\section{Reinforcement Learning Approach}
\label{sec:rl-approach}

TODO: Belohnungsfunktion, DQN-Agent (Trainingsschleife, Replay-Buffer, Zielnetzwerk, $\epsilon$-greedy).

To use reinforcement learning techniques for modeling the \ac{drausp}, first of all we need to define
a \ac{mdp} that represents the problem. We define a state space,an action space and a transition function: 

\textbf{State space.} A state
\[
  s = (t, \hat C_1, \dots, \hat C_K, r, q_1, \dots, q_K) \in \mathcal{S} \subseteq \mathbb{R}^{2K+2}
\]
consists of the current decision epoch $t \in \{1,\dots,T_d\}$, the remaining capacities
$\hat C_k \in \{0,\dots,C_k^{(0)}\}$ of each resource $k=1,\dots,K$, and the $t$-th request $[r,q]$
currently up for decision, with reward $r \in \mathbb{R}_{\ge 0}$ and (zero-padded) demand vector
$q \in \mathbb{N}_0^K$:
\begin{equation}
  \label{eq:padding}
  q = (\underbrace{q_1, \dots, q_\ell}_{\text{demanded capacities}},\; \underbrace{0, \dots, 0}_{K-\ell \text{ padding}}\,).
\end{equation}
The initial state is $s_1 = (1, C_1^{(0)}, \dots, C_K^{(0)}, r_1, q_1)$.

\textbf{Action space.} $\mathcal{A} = \{0,1,\dots,K\}$: action $a=0$ rejects the current request; action $a \in \{1,\dots,K\}$ accepts it and schedules it starting at slot $a$.
Assume that the current request has length $\ell \in \{1,\dots,K\}$ and demand vector $q = (q_1,\dots,q_\ell,0,\dots,0)$ (zero-padded to length $K$). Then the action $a$ is feasible if and only if 
\begin{itemize}
 \item $a+\ell-1 \le K$, we do not exceed the maximum length of resources
 \item $\hat C_{a+i-1} \ge q_i$ for all $i=1,\dots,\ell$, the capacity of all affected slots is sufficient to accept the request
\end{itemize} 
Note that for a given state not all actions are feasible: We define the feasible action set $\mathcal{A}_s^{\text{valid}} \subseteq \mathcal{A}$ as the set of actions that do not exceed the remaining capacity of any affected slot and do not exceed the horizon $K$.
Nevertheless the set of valid actions $\mathcal{A}_s^{\text{valid}}$ is never empty, since $0 \in \mathcal{A}_s^{\text{valid}}$ for all states $s$.
We will use the set of valid actions later on.

\textbf{Transition function.} $\mathcal{T}: \mathcal{S} \times \mathcal{A} \to \mathcal{S}$
deterministically updates the remaining capacities according to the chosen action, appends the next request
$[r_{t+1}, q_{t+1}]$ drawn from the instance's request sequence, and advances time, $t \mapsto t+1$. 
The new capacities are given by the following equation:
\[
  \hat C_{a+i-1} \;\leftarrow\; \hat C_{a+i-1} - q_i, \qquad i = 1,\dots,\ell.
\]
Note that $\hat{C_k} \geq 0$ for all valid actions $a \in \mathcal{A}_x^{\text{valid}}$ and all $k=1,\dots,K$.

Figure~\ref{fig:state-transition} illustrates the effect of the transition function $\mathcal{T}$ on a concrete example with
$K=5$, request length $\ell=3$, action $a=2$ and applies the update rule.

\begin{figure}[htbp]
  \centering
  \begin{tikzpicture}[font=\small,
      slot/.style={draw, minimum width=11mm, minimum height=11mm},
      occ/.style ={draw, fill=orange!25, minimum width=11mm, minimum height=11mm},
      >={Latex[length=2mm]}
    ]
    % Before-row: state s at epoch t
    \node[anchor=east] at (-0.5,0)  {before, $s_t$:};
    \node[slot] (b1) at (0,0)   {$\hat C_1$};
    \node[occ]  (b2) at (1.3,0) {$\hat C_2$};
    \node[occ]  (b3) at (2.6,0) {$\hat C_3$};
    \node[occ]  (b4) at (3.9,0) {$\hat C_4$};
    \node[slot] (b5) at (5.2,0) {$\hat C_5$};

    \node[below=1pt, gray] at (b1.south) {\scriptsize $k=1$};
    \node[below=1pt, gray] at (b2.south) {\scriptsize $k=2$};
    \node[below=1pt, gray] at (b3.south) {\scriptsize $k=3$};
    \node[below=1pt, gray] at (b4.south) {\scriptsize $k=4$};
    \node[below=1pt, gray] at (b5.south) {\scriptsize $k=5$};

    \draw[decorate, decoration={brace, amplitude=5pt}] (1.85+0.05,0.65) -- (3.55-0.05,0.65)
      node[midway, above=5pt] {action $a=2$, request length $\ell=3$};

    % After-row: state s_{t+1} after applying T
    \node[anchor=east] at (-0.5,-3) {after, $s_{t+1}$:};
    \node[slot] (a1) at (0,-3)   {$\hat C_1$};
    \node[occ]  (a2) at (1.3,-3) {$\hat C_2{-}q_1$};
    \node[occ]  (a3) at (2.6,-3) {$\hat C_3{-}q_2$};
    \node[occ]  (a4) at (3.9,-3) {$\hat C_4{-}q_3$};
    \node[slot] (a5) at (5.2,-3) {$\hat C_5$};

    \draw[->, dashed, gray] (b1) -- (a1);
    \draw[->] (b2) -- node[right, xshift=1pt] {$-q_1$} (a2);
    \draw[->] (b3) -- node[right, xshift=1pt] {$-q_2$} (a3);
    \draw[->] (b4) -- node[right, xshift=1pt] {$-q_3$} (a4);
    \draw[->, dashed, gray] (b5) -- (a5);
  \end{tikzpicture}
  \caption{Effect of the transition function $\mathcal{T}$ ($K=5$ slots, request length $\ell=3$, action $a=2$): the
    remaining capacities of the three occupied slots each decrease by the corresponding entry of
    the demand vector $q$, while the capacities of slots outside the accepted block ($k=1,5$)
    are unaffected.}
  \label{fig:state-transition}
\end{figure}

\textbf{Reward.} The reward function $\mathcal{R}: \mathcal{S} \times \mathcal{A} \to \mathbb{R}$ is defined as follows:
\[
  \mathcal{R}(s,a) =
  \begin{cases}
    r, & \text{if } a \in \mathcal{A}_s^{\text{valid}} \setminus \{0\} \\
    0, & \text{if } a = 0 \\
    -100, & \text{if } a \notin \mathcal{A}_s^{\text{valid}}
  \end{cases}
\]
Selecting an infeasible action, terminates the episode and yields a large negative reward of $-100$ to discourage the agent from taking invalid actions.

\textbf{Objective.} 
Given a MDP the goal is to find a policy $\pi$: $S$ $\to$ $A$ that maximizes the expected cumulative reward over the decision horizon. In our case, the objective is to maximize the total revenue over the decision horizon by learning an acceptance-and-scheduling policy that maps states to actions.
\begin{equation*}
      \pi^* =\arg\max_{\pi} \mathbb{E}_{\pi}\Bigg[\sum_{t=1}^{\infty}\gamma^{t} r_t\Bigg]
  \end{equation*}
Solving this optimization directly is a hard task because we have to find the optimal policy $\pi^*$ over the whole state space and the policy $\pi$ is itself part of the expectation. Instead, we will use Q-learning to learn the optimal action-value function $Q^*(s,a)$, which allows us to derive the optimal policy $\pi^*$ by acting greedily with respect to it.

\textbf{(Deep) Q-Learning.} 
In Q-learning, we learn the optimal action-value function
\[
  Q^*(s,a) = \max_\pi \; \mathbb{E}\Big[\, \sum_{t'=t}^{T_d} \gamma^{t'-t} \cdot \mathcal{R}(s_{t'}, a_{t'}) \;\Big|\; s_t = s,\, a_t = a \,\Big].
\]
Once $Q^*$ is known (or
approximated by a trained network), the optimal policy follows directly by acting greedily with
respect to it:
\[
  \pi^*(s) = \arg\max_{a \in \mathcal{A}} Q^*(s,a).
\]
Remark that Q-Learning suffers from the curse of dimensionality, since the state space grows exponentially with the number of resources $K$ and the number of requests $T_d$. Therefore, we will use a Deep Q-Network (DQN) 
$$
Q(\cdot \vert \theta_{DQN}): \quad s \xrightarrow{\text{DQN}}   \begin{pmatrix} Q(s, a_1) \\ Q(s, a_2) \\ \vdots \\ Q(s, a_N) \end{pmatrix}
$$
to approximate the action-value function $Q^*(s,a)$.

We use standard DQN architecture as a baseline which tries to learn the action-value function $Q(s,a)$ without any structural knowledge about the problem. 

\textbf{Comparison with the state/action space formulation in the DRAUSP paper.}
\textcite{drausp_lion18} also formulate \acs{drausp} as an \acs{mdp}, but with a considerably
leaner state space. Their state is simply the capacity vector
$$s = (\hat C_1,\dots,\hat C_K) \in S = \{0,\dots,C_1\}\times\dots\times\{0,\dots,C_K\}$$ the
time index $t$ is not part of the state vector at all, but only a separate, downward-counting step
index of their recursion, indicating how many requests are still to arrive. In particular, the
current request $[r,q]$ is \emph{not} contained in the state. Our state $s$ (Equation~\eqref{eq:state}), by contrast, contains $r$ and $q$
explicitly, while our action is reduced to the lean start slot $a \in \{0,\dots,K\}$.
Table~\ref{tab:mdp-comparison} summarizes this structural comparison; \emph{how} the two
formulations are subsequently solved (exact dynamic programming vs.\ function approximation via
\acs{dqn}) is discussed separately in the section on the DQN agent.

\begin{table}[htbp]
  \centering
  \begin{tabular}{p{0.20\textwidth}p{0.36\textwidth}p{0.36\textwidth}}
    \toprule
     & \textbf{Paper} & \textbf{Our Approach} \\
    \midrule
    State    & $s_t = (\hat C_1,\dots,\hat C_K)$, $t$ a separate step index
             & $s = [t,\hat C_1,\dots,\hat C_K,r,q_1,\dots,q_K]$ \\
    Request  & embedded in the action set $A_s$
             & embedded explicitly in the state \\
    Action   & $a=[a_r,a_q]$: full, request-dependent $K$-vector per start position
             & $a \in \{0,\dots,K\}$: only the start slot \\
    \bottomrule
  \end{tabular}
  \caption{Comparison of the state/action space formulation in the DRAUSP paper
    \parencite{drausp_lion18} and in our approach.}
  \label{tab:mdp-comparison}
\end{table}

The paper's state space is insufficient for a Q-learning approach, because a Q-function $Q(s,a)$ is
supposed to return one fixed value per state $s$ for every action $a$ -- but here, how good an
action is depends crucially on which request is currently at hand.

At every decision point, a Q-learning agent should be able to read off, directly from the state,
how good each action is. This requires the state to include the request itself, so that the agent
knows both the reward and the resource demand it is deciding about.


\subsection{Zustands- und Aktionsraum}
\label{subsec:state-action}

Das DRAUSP verwaltet $K$ fest angeordnete Kapazitätsslots (Ressourcen) $k=1,\dots,K$ mit
Anfangskapazitäten $C_k^{(0)}$ (z.\,B.\ $C_k^{(0)}=20$ für alle $k$), die über einen
Planungshorizont von $T_d$ Entscheidungsepochen \emph{nie} wieder aufgefüllt werden. Der Zustand
zu Beginn von Epoche $t$ lässt sich als
\begin{equation}
  \label{eq:state}
  s = \big[\, t,\; \hat C_1, \dots, \hat C_K,\; r,\; q_1, \dots, q_K \,\big] \in \mathbb{R}^{2K+2}
\end{equation}
darstellen, wobei $t$ die verstrichene Zeit, $\hat C_k$ die verbleibende Restkapazität von Slot
$k$, $r$ die Belohnung der in Epoche $t$ eintreffenden Anfrage und $q_k$ deren Bedarfsvektor
bezeichnet (\texttt{gymnasium\_env.py}, Klasse \texttt{DrauspEnv}).

\textbf{Anfragen haben variable Länge, der Zustand nicht.} Eine Anfrage belegt bei Annahme keinen
einzelnen Slot, sondern einen \emph{zusammenhängenden Block} von $\ell \in \{1,\dots,K\}$ Slots,
und benötigt dabei für den $i$-ten Slot dieses Blocks $q_i$ Kapazitätseinheiten,
$i = 1,\dots,\ell$. Die Blocklänge $\ell$ variiert von Anfrage zu Anfrage -- im
Instanzgenerator (\texttt{instance\_generator.py}) wird sie gleichverteilt aus
$\{1,\dots,K\}$ gezogen; in Benchmark-Instanzen kann sie auch fix vorgegeben sein
(\texttt{fixed\_request\_length}). Da aber \emph{ein und dasselbe} Netzwerk Zustände zu jeder
Anfrage verarbeiten muss, wird der Bedarfsvektor unabhängig von $\ell$ stets auf die feste Länge
$K$ gebracht, indem er mit Nullen aufgefüllt wird:
\begin{equation}
  \label{eq:padding}
  q = (\underbrace{q_1, \dots, q_\ell}_{\text{tatsächlicher Bedarf}},\; \underbrace{0, \dots, 0}_{K-\ell \text{ Nullen}}\,).
\end{equation}
Entscheidend dabei: Der Index $k$ in $q_k$ bezeichnet \emph{keine} absolute Slotposition im
Zeitplan, sondern die Position \emph{relativ zum (noch nicht feststehenden) Startslot} der
Anfrage. Wo genau dieser Block im Zeitplan liegt, legt erst die Aktion fest.

\textbf{Aktion.} Die Aktion $a \in \{0,1,\dots,K\}$ wählt entweder Ablehnung ($a=0$, Reward $0$)
oder den Startslot $a$, ab dem die Anfrage eingeplant wird. Bei Annahme belegt sie die Slots
$a, a+1, \dots, a+\ell-1$ und reduziert deren Restkapazität um $q_1,\dots,q_\ell$:
\[
  \hat C_{a+i-1} \;\leftarrow\; \hat C_{a+i-1} - q_i, \qquad i = 1,\dots,\ell.
\]
Eine Aktion ist nur gültig, wenn der Block innerhalb des Horizonts liegt ($a+\ell-1 \le K$) und an
jeder betroffenen Slotposition ausreichend Restkapazität vorhanden ist ($\hat C_{a+i-1} \ge q_i$
für alle $i$); ungültige Aktionen werden mit Reward $-20$ bestraft und beenden die Episode.
Abbildung~\ref{fig:state-padding} veranschaulicht dies für $K=5$, eine Anfrage der Länge
$\ell=3$ und Startslot $a=2$.

\begin{figure}[htbp]
  \centering
  \begin{tikzpicture}[font=\small,
      slot/.style={draw, minimum width=11mm, minimum height=11mm},
      occ/.style ={draw, fill=orange!25, minimum width=11mm, minimum height=11mm},
      >={Latex[length=2mm]}
    ]
    % Slots 1..5, Slots 2-4 durch die Anfrage belegt (a=2, l=3)
    \node[slot] (s1) at (0,0) {$\hat C_1$};
    \node[occ]  (s2) at (1.3,0) {$\hat C_2$};
    \node[occ]  (s3) at (2.6,0) {$\hat C_3$};
    \node[occ]  (s4) at (3.9,0) {$\hat C_4$};
    \node[slot] (s5) at (5.2,0) {$\hat C_5$};

    \node[below=1pt, gray] at (s1.south) {\scriptsize $k=1$};
    \node[below=1pt, gray] at (s2.south) {\scriptsize $k=2$};
    \node[below=1pt, gray] at (s3.south) {\scriptsize $k=3$};
    \node[below=1pt, gray] at (s4.south) {\scriptsize $k=4$};
    \node[below=1pt, gray] at (s5.south) {\scriptsize $k=5$};

    % Block-Klammer über den belegten Slots
    \draw[decorate, decoration={brace, amplitude=5pt}] (1.85+0.05,0.65) -- (3.55-0.05,0.65)
      node[midway, above=5pt] {Block $\ell=3$, Startslot $a=2$};

    % q_i Zuordnung -- Pfeile starten unterhalb der k-Beschriftung, damit nichts überlappt
    \draw[->] (s2.south) ++(0,-0.5) -- ++(0,-0.9) node[below] {$q_1$};
    \draw[->] (s3.south) ++(0,-0.5) -- ++(0,-0.9) node[below] {$q_2$};
    \draw[->] (s4.south) ++(0,-0.5) -- ++(0,-0.9) node[below] {$q_3$};

    % Roher, gepolsterter Bedarfsvektor der Anfrage
    \node at (2.6,-3.0)
      {$q = \big(\, \underbrace{q_1,\,q_2,\,q_3}_{\text{Bedarf, relativ zu } a},\;
                    \underbrace{0,\,0}_{\text{Padding}} \,\big) \in \mathbb{R}^5$};
  \end{tikzpicture}
  \caption{Zustandsaufbau bei $K=5$ Slots und einer Anfrage der Länge $\ell=3$: Der Bedarfsvektor
    $q$ ist relativ zum (von der Aktion bestimmten) Startslot indiziert und mit $K-\ell=2$ Nullen
    auf feste Länge~$K$ aufgefüllt. Erst die gewählte Aktion $a$ ordnet $q_1,q_2,q_3$ den
    absoluten Slots $\hat C_2,\hat C_3,\hat C_4$ zu.}
  \label{fig:state-padding}
\end{figure}

\textbf{Konsequenz für die Monotonie.} Die hinteren $K-\ell$ Einträge von $q$ sind stets $0$ --
der kleinste Wert im kalibrierten Wertebereich von $q_k$. Das ist inhaltlich konsistent mit der
in Abschnitt~\ref{subsec:monotonicity} postulierten DECREASING-Monotonie: Ein Padding-Eintrag
bedeutet "kein Bedarf an dieser relativen Position" und wirkt im kalibrierten Modell wie eine
Anfrage ohne jede Einschränkung an diesem Slot -- was korrekt ist, da dieser Slot vom Block der
aktuellen Anfrage gar nicht betroffen ist.
